\cotefile{162-5}
\scanpage{28}
7.\quad \uline{Coefficients} \inserted{\uline{sur catégorie formelle}} Soit $M$ un $\underline O_{X_0}$-module. La notion de ``donnée de descente'' sur $M$ relativement à $\mathcal C$ [ou ``$\mathcal C$-stratification'', ou ``opération de $\mathcal C$ sur $M$''] est claire. On constate alors \unsure{qu'une|que une} telle structure permet de mettre sur
\[
\bigwedge^*\Omega\otimes_{\underline O_{X_0}}M
\]
un opérateur différentiel $\delta_M$, tel que
\[
\begin{array}{l@{\qquad}l}
(*)&\delta_M(fx)=\delta(f)x\deleted{\unsure{\delta_M(x)}}+(-1)^{\deg F}f\,\delta_M(x).
\end{array}
\]
[Il suffit $=$ d'avoir une donnée de recollement d'ordre 1]

Un tel opérateur est connu quand on connait
\[
M\xrightarrow{\delta_M^1}\Omega\otimes M
\]
satisfaisant $(*)$ \deleted{\unsure{$\delta_M^1($}} pour $f\in\Gamma\underline O_{X_0}$ (donc $\deg f=0$).

\unsure{Le|La} fait que l'opérateur déduit \inserted{sur $\bigwedge^*\Omega\otimes M$ cf plus bas)} soit de carré nul provient de la condition de descente. [il suffit qu'elle soit vraie à l'ordre 2 \ldots??].

8.\quad \uline{Coefficients relatifs à une Algèbre de De Rham}

Plus gén, si $(\bigwedge^*\Omega,\delta)$ est une alg de De Rham sur $X_0$, on \deleted{\unsure{peut}} regarde les modules $M$ avec une \uline{connexion} relativement à celle-ci, i.e. un morphisme
\[
M\overset{\delta_M}{\mathbin{-}}\ \Omega\otimes M
\]
satisfaisant
\[
\delta_M(fx)=\delta(f)\otimes x+f\,\delta_M(x).
\]

\margin{27}
